\chapter{Results and Discussion}\label{ch:results}

This chapter brings together what the finished system does, what the measurements showed,
and how far the original objectives were met. Wherever a number appears, it comes from the
runs described in Chapters~\ref{ch:perf} and~\ref{ch:test}; nothing here is estimated.

\section{Implemented System at a Glance}
The delivered system has three running parts: a MySQL database, a REST backend and a
server-rendered web portal. Table~\ref{tab:delivered} summarises the size of each part.

\begin{table}[H]
\centering\small
\caption{Size of the delivered system}\label{tab:delivered}
\begin{tabularx}{\linewidth}{@{}p{5.4cm}X@{}}
\toprule \textbf{Part} & \textbf{Measured size}\\ \midrule
Database & 63 tables in 13 functional modules; 91 foreign keys; 25 ENUM columns; 43 unique keys; 6 views; 5 triggers; 3 stored procedures; 25 added indexes\\
Backend (Java) & 152 source files, about 11\,300 lines of main code and 4\,400 lines of test code; 225 operations on 151 URL paths\\
Portal (Java and templates) & 18 controllers (about 4\,900 lines), 70 page templates (about 8\,300 lines)\\
Security data & 7 default roles and 62 permissions, loaded from the database on every request\\
Automated tests & 185 integration tests in 15 classes, 5 of them concurrency tests; a 42-check SQL suite\\
\bottomrule
\end{tabularx}
\end{table}

\section{Implemented Features}
Table~\ref{tab:features} maps each functional area to what a user can actually do in the
portal. Every row was exercised through the browser with the demonstration data.

\begin{table}[H]
\centering\small
\caption{Features available in the portal}\label{tab:features}
\begin{tabularx}{\linewidth}{@{}p{3.6cm}X@{}}
\toprule \textbf{Area} & \textbf{What can be done}\\ \midrule
Citizens & Register with address; search and filter; edit; blacklist with a reason; view addresses and linked records\\
Applications & Create for a service type and office; follow the status workflow; attach documents; book appointment slots; view history\\
Vehicles & Register; view current owner and ownership history; request, approve or reject an ownership transfer; view compliance records\\
Compliance & Fitness certificates, insurance, pollution certificates and road tax, with an \emph{expiring soon} view\\
Licences & Learner and driving licences, driving tests and attempts, driving schools\\
Permits and routes & Commercial permits and the transport routes they cover\\
Violations and challans & Record violations, issue e-challans, follow challan status, file complaints and legal appeals\\
Payments & Pay fees and challans through attempts and receipts; replaying the same gateway reference is safe; refunds\\
Administration & Users, roles and permissions, employees, offices and counters, appointment slots, fee structures, reference data, notifications and the audit trail\\
Citizen self-service & A reduced five-item menu for a signed-in citizen: dashboard, applications, complaints, appeals and notifications, limited by the permissions of the citizen role\\
\bottomrule
\end{tabularx}
\end{table}

\section{Screens of the Running System}
Six screens are shown here; the rest of the portal follows the same layout of a left menu,
a title, and cards or tables. All screens were captured from the running application with
the demonstration data.

Figure~\ref{fig:scrdash} is the administrator dashboard, the first page after sign-in. It
shows counts that are read from the backend summary endpoint and shortcuts to the most
common work.
\fig{03_dashboard.png}{Administrator dashboard}{fig:scrdash}{0.58}

Figure~\ref{fig:scrapp} shows one application with its current status, the allowed next
steps, attached documents and booked appointments. Only transitions permitted by the
state diagram of Chapter~\ref{ch:design} are offered as buttons.
\fig{09_application_detail.png}{Application detail with status workflow}{fig:scrapp}{0.58}

Figure~\ref{fig:scrveh} is the vehicle page. The current owner is shown separately from the
ownership history, and the compliance documents are listed with their validity dates.
\fig{12_vehicle_detail.png}{Vehicle detail with current owner and history}{fig:scrveh}{0.68}

Figure~\ref{fig:scrrole} shows a role with the permissions assigned to it. An administrator
changes access here without any change to the program, which is the point of keeping the
permissions in the database.
\fig{38_admin_role_detail.png}{Role detail with its permissions}{fig:scrrole}{0.58}

Figure~\ref{fig:citizenmenu} is what a citizen sees. The menu has only five items (dashboard,
applications, complaints, appeals, notifications) and the administrator pages are refused, as tested in Section~\ref{sec:crawl}.
\fig{44_citizen_dashboard.png}{Citizen dashboard with the reduced menu}{fig:citizenmenu}{0.42}

Finally, Figure~\ref{fig:scrval} shows input checking. A future date of birth, a national
identity number that is too short and a phone number with letters are all rejected by the
backend. The portal keeps what the user typed, outlines each wrong field and prints the
message under it. This screen also records one of the defects found late in the work
(D-10 in Table~\ref{tab:defects}): in the first version the typed values were lost.
\fig{06b_citizen_validation_errors.png}{Field-level validation messages on the citizen form}{fig:scrval}{0.62}

\section{Functional Results}
All 24 functional requirements of Chapter~\ref{ch:srs} are traced to at least one
passing test in Table~\ref{tab:tc}. The headline results are as follows.
\begin{itemize}
\item \textbf{185 of 185} backend integration tests pass against a real MySQL 8.0 server.
\item \textbf{5 of 5} concurrency tests pass: two users cannot take the last appointment slot, a payment replayed with the same gateway reference is recorded once, and parallel status changes or transfer approvals leave exactly one winner.
\item \textbf{91 of 91} foreign keys have no orphan rows after the seed is loaded.
\item The browser crawl ended with \textbf{no server errors}; the only responses that were not successful were the 405 answers expected when a post-only address is opened in a browser (Table~\ref{tab:crawl}).
\item The SQL suite passes \textbf{39 of 42} checks; the three that fail are explained in Section~\ref{sec:sqltests}; they come from an out-of-date expected value and from the way the script is run (defect D-09), not from the running system.
\end{itemize}

\section{Performance Results}
Table~\ref{tab:perfsummary} repeats the key numbers of Chapter~\ref{ch:perf} against the
target set in Table~\ref{tab:kpi}.

\begin{table}[H]
\centering\small
\caption{Performance summary}\label{tab:perfsummary}
\begin{tabularx}{\linewidth}{@{}p{5.2cm}p{3.6cm}X@{}}
\toprule \textbf{Measure} & \textbf{Result} & \textbf{Comment}\\ \midrule
Mean latency, 1 user & 22.6~ms & Read endpoints, one machine\\
Mean latency, 20 users & 46.2~ms & 95th percentile 158~ms\\
Throughput, 20 users & 380 req/s & No errors in 2\,160 requests\\
Sign-in & 445~ms & Deliberately slow (BCrypt cost 12)\\
Lookup on 500\,000 rows & 405.9~ms $\rightarrow$ 1.4~ms & About 290 times faster with the index\\
Backend memory & about 110~MB & Idle after the tests\\
\bottomrule
\end{tabularx}
\end{table}

\section{Comparison with the Original Objectives}
Table~\ref{tab:objectives} compares each objective of Chapter~\ref{ch:intro} with the
outcome. We have marked an objective as met only where there is evidence for it in this
report.

\begin{table}[H]
\centering\small
\caption{Objectives against outcomes}\label{tab:objectives}
\begin{tabularx}{\linewidth}{@{}p{4.4cm}Xp{1.8cm}@{}}
\toprule \textbf{Objective} & \textbf{Evidence} & \textbf{Status}\\ \midrule
A normalised schema for the main RTO processes & 63 tables, 91 foreign keys, no orphan rows; Chapter~\ref{ch:design} & Met\\
Rules enforced by the database, not only by code & 19 applied CHECK rules, unique keys, generated columns, 5 triggers; Section~\ref{sec:schemaaudit} & Met, with D-08\\
Correct behaviour under concurrent use & Ordered locks, guarded updates; 5 concurrency tests; Section~\ref{sec:conc} & Met\\
Role-based access with auditability & 7 roles, 62 permissions, audit trail in the same transaction; Section~\ref{sec:rbactest} & Met\\
A usable web portal for staff and citizens & 70 templates, crawl of about 820 requests; Section~\ref{sec:crawl} & Met\\
Measured, not assumed, performance & Load test and index experiment; Chapter~\ref{ch:perf} & Met\\
Use of the course subjects in one system & Chapter~\ref{ch:subjects} & Met\\
\bottomrule
\end{tabularx}
\end{table}

\section{Discussion}
\paragraph{What worked well.} Putting the rules in more than one layer paid off. The
validation of the API, the transaction logic of the services and the constraints of the
database each caught different mistakes during testing, and none of the three alone would
have been enough. Reading permissions from the database made the role screen real: it
changes access immediately and it appears in the audit trail.

\paragraph{What was harder than expected.} Concurrency was the most subtle part. A
check-then-write sequence looks correct in a single-user test and fails only when two
requests interleave; the cure was to lock rows in a fixed order and to make the write itself
conditional. The other surprise was how many faults sat outside the main logic: error pages,
seed data and form handling accounted for most of the defects in Table~\ref{tab:defects}.

\paragraph{What the numbers do and do not show.} The load test used one laptop and read
requests only, so the figures describe this machine and this kind of work. They support the
claim that the system is comfortable at the scale of one office. They do not support any
claim about a national deployment.

\section{Practical Applicability}
In its present form the system suits a pilot at a single regional office or a teaching
environment, where staff handle citizens, applications, vehicles, licences and fines through
one consistent interface. The structure, with a clear API between the portal and the
data, would also let a mobile application or another office system reuse the same backend.
Connection to real payment gateways, SMS providers and national registers is outside the
scope of this work and is listed in Chapter~\ref{ch:conclusion} as future scope.
